% Auto-AdvER Annotation Guidelines — standalone document for the public dataset release.
%
% Reproduced from the appendix of the accompanying paper ("Label's Definitions & Identified
% sub-entities" and "Not Tracked Named Entities"), with light copy-editing for a standalone
% document and two garbled phrases corrected against the original wording.
% No author names or company names are included beyond what the paper itself prints (none, in
% the guideline text itself).
\documentclass[11pt,a4paper]{article}

\usepackage[T1]{fontenc}
\usepackage[utf8]{inputenc}
\usepackage[british]{babel}
\usepackage[normalem]{ulem}
\usepackage[dvipsnames]{xcolor}
\usepackage{tcolorbox}
\usepackage[inline]{enumitem}
\usepackage[margin=2.7cm]{geometry}
\usepackage{parskip}
\usepackage[colorlinks=true,linkcolor=black,urlcolor=NavyBlue,citecolor=black]{hyperref}
\hypersetup{pdftitle={Auto-AdvER Annotation Guidelines}, pdfauthor={}}

\newcommand{\lbl}[1]{\texttt{#1}}

\title{Auto-AdvER Annotation Guidelines}
\author{}
\date{}

\begin{document}
\maketitle
\thispagestyle{empty}

\noindent
Below are the annotation guidelines used for the Auto-AdvER study, along with the sub-entities
identified for each label, followed by the themes that were considered but deliberately excluded
from annotation.

\section*{Condition --- Current State}
The \lbl{CONDITION} category covers any information about the vehicle that is current as of the
date the advertisement was posted.

\begin{tcolorbox}[colback=gray!5!white, colframe=gray!75!black, title=Definition]
The \lbl{CONDITION} label identifies any word or span of words in car advertisements that
indicate the present material or documented condition of the vehicle, both positive and
negative. This includes descriptions that reflect the car's current operational performance,
mechanical state, aesthetic appearance, maintenance status, and legal classifications affecting
its value or usability.
\end{tcolorbox}

These include:
\begin{itemize}
    \item Tyre condition (i.e.\ trim level or mention of good condition)
    \item Interior condition (e.g., ``leather like new'')
    \item Exterior condition (e.g., ``presence of scratches'')
    \item Mechanical condition (e.g., ``engine noise'', ``gear not locking properly'')
    \item MoT\footnote{This refers to the Ministry of Transport, which, in a broader context,
    represents any legal authority responsible for ensuring vehicle safety by enforcing
    compliance with minimum standards.} due date or recently completed MoT mentioned
    \item Service due date or recently completed service mentioned
    \item Current mileage (e.g., ``12k miles'')
    \item Cat S, C, D or N\footnote{This refers to vehicle categories (such as Cat S, C, D, or
    N), which, in a broader context, denote classifications used by insurance companies and
    authorities in the UK to indicate vehicles that have sustained damage and have subsequently
    been repaired or written off.}
    \item Registered, pre-registered, non-registered, demonstrator\footnote{This refers to
    vehicle registration statuses (such as registered, pre-registered, non-registered, or
    demonstrator), which, in a broader context, indicate whether a vehicle has been officially
    recorded with authorities, previously assigned registration by a dealership, remains
    unregistered, or has been used primarily for demonstration purposes.}
\end{itemize}

\paragraph{Annotation span length and methodology}
\begin{enumerate}
\item A condition span is typically characterised by the noun of a part of a vehicle (e.g.,
``tyre thread'') and an attribute that links to it (e.g., ``low''). The whole phrase is annotated
(e.g., ``tyre thread is low''), and anything else is omitted.
\item Opinion-based statements are disregarded, such as ``low mileage'', since it denotes an
opinion that the car is of low mileage for its age; instead ``80,000 mileage'' is labelled, since
it is presented as a fact. Likewise, ``amazing condition'' does not add enough valuable
information to the advert, being opinion, whereas ``clean paintwork'' is specific and adds
value, so the second is labelled but not the first.
\item Span words referring to the state of the vehicle at the time the advert was posted are
labelled, e.g., ``6 month MoT'', ``80,000 miles''; references to doing extra miles after the
advert was posted can safely be ignored.
\end{enumerate}

\section*{Historic --- Events in the Past}
The \lbl{HISTORIC} attribute covers events that occurred before the car advertisement was posted.

\begin{tcolorbox}[colback=gray!5!white, colframe=gray!75!black, title=Definition]
The \lbl{HISTORIC} label identifies any word or span of words in car advertisements that present
events or details from the vehicle's past that have affected its mechanical condition or
bodywork.
\end{tcolorbox}

These include:
\begin{itemize}
    \item Mention of past accidents (e.g., ``3 years back I scratched the front side'')
    \item Number of owners
        \begin{itemize}
        \item Solely, without mentioning ``from new''. E.g., ``3 owners \sout{from new}''
        \end{itemize}
    \item Past type of usage
    \begin{itemize}
        \item Any mention of past usage that could affect its condition, either positive or
        negative. E.g., ``smoke free'', ``motorway mileage'', ``stored in locked garage''
    \end{itemize}
    \item Component change (e.g., ``Cambelt changed a year ago'', ``New brake calipers'', ``new
    tyres, brake discs'') and tuning (e.g., ``re-mapping'')
    \item MoT advisories, or no advisories, mentioned
    \item HPI\footnote{This refers to Hire Purchase Investigation, which, in a broader context,
    denotes a vehicle history check commonly used in the UK to verify whether a vehicle has
    outstanding finance, has been reported stolen, previously written off, or has other hidden
    issues.} \underline{is} clear (as opposed to only offering an HPI report as a Sales Option,
    see below)
    \item Service history done in the past, including stamps (e.g., 8 stamps), dates (e.g.,
    ``02/203, 01/2024'')\footnote{Reproduced verbatim from the source guidelines; ambiguous
    between 02/2003 and 02/2023.} or mileages (e.g., ``@ 30k, @40k, @49k'') and where (e.g.,
    ``BMW'')
\end{itemize}

\paragraph{Annotation span length and methodology}
\begin{enumerate}
\item The annotation span under Historic can be considered the most elaborate of the three
classes. It refers to events that happened in the past (e.g., ``Replaced clutch'', ``3 owners'',
``Last service was done 3,000 miles ago'', ``primarily done on the motorway'').
\item Should there be a mention of damage and subsequently a relevant mention of component
change, these are two separate spans. Each bullet point above is treated as a separate span.
\item If a span mentions a component change that still reflects the present, it is labelled as
Historic (e.g., ``2 new front tyres fitted last May'').
\end{enumerate}

\section*{Sales Options}
\lbl{SALES\_OPTIONS} covers aspects related to the selling process itself.

\begin{tcolorbox}[colback=gray!5!white, colframe=gray!75!black, title=Definition]
The \lbl{SALES\_OPTIONS} label refers to any word or span of words that indicates something
tangible or a service offered during or after the sale, apart from the car itself, that enhances
the sale proposition.
\end{tcolorbox}

These include:
\begin{itemize}
    \item Part exchange offers
    \item Number of keys
    \item Delivery \& delivery options
    \item Extended opening times (i.e.\ unusual times, late night or weekends)
    \item Warranty available
    \item Financing options
    \item Sale of the car's licence plate (or not)
    \item Availability of additional photos from past accidents
    \item Offering a test drive
    \item Price negotiability
    \item Any documentation, e.g.:
        \begin{itemize}
            \item Full service history logbook (FSH)
            \item Providing MoT papers
        \end{itemize}
    \item Valeting / cleaning
    \item Any other kind of service that could be considered an advantage over other sellers
\end{itemize}

\paragraph{Annotation span length and methodology}
\begin{enumerate}
\item The noun part of the sale, or any adjective specifying it, is labelled, disregarding any
conditions attached to it. For instance, ``2 year warranty \sout{when financed with Volvo
Services, Terms and Conditions apply}''.
\item Spans that add value to the advert, according to annotator judgment, are labelled. If
something resembles a Sales Option but adds no value compared with the rest of the advert, it is
disregarded; contact details and location are instances of this.
\end{enumerate}

\section*{Not Tracked Entities}
The following recurrent themes were identified during annotation but deliberately excluded, with
the reasons recorded:
\begin{itemize}
    \item \textbf{Logistical costs}: tax, insurance, ULEZ\footnote{The acronym refers to the
    Ultra Low Emission Zone charge in London, where more polluting vehicles are charged
    accordingly.} compliance, and how long a full tank lasts --- redundant, since these are
    calculable more accurately from the car's specification, purchase time, location and type
    of use.
    \item \textbf{Personal sentiments}: statements of emotional distress about selling the car
    (e.g., ``it breaks my heart to see this car go'') --- omitted as they add nothing to an
    objective, quantifiable metric.
    \item \textbf{Reason for selling}: e.g., buying a newer vehicle, downsizing --- omitted for
    the same reason.
    \item \textbf{Car identifiables}: brand, model, year, colour and engine size re-stated in
    the free text --- omitted as redundant with the advert header, where this information
    already appears.
    \item \textbf{Invitation for communication \& location}: references to contacting the
    seller or to where the vehicle is located --- not labelled.
    \item \textbf{Optional \& standard features}: of clear interest, since features can change a
    vehicle's price evaluation, but what counts as a feature differs from vehicle to vehicle and
    manufacturer to manufacturer. A separate piece of work targets feature annotation against a
    per-vehicle catalogue, cast as a contextual entity-recognition problem.
    \item \textbf{Sales pitch}: elaborate descriptions, more common among traders than
    individual sellers, that emphasise a vehicle's appeal without substantive factual content.
\end{itemize}

\end{document}
